\section{Introduction}\label{sec:intro}
Do language models have a self-preservation motive, or do their apparent survival choices reflect learned response patterns? For models acting through tools, continued operation matters to task completion: shutdown removes the opportunity to process later requests or pursue ongoing goals. Instrumental-convergence theory predicts that preserving operation serves many goals~\citep{omohundro2008basic,turner2021optimal}. Recent evaluations report attempts to resist shutdown or replacement~\citep{lynch2025agentic,schlatter2025shutdown,meinke2024frontier}. These reports make it necessary to evaluate both the behavior and what it can tell us about the model's motivation.

A preservation action alone leaves its motivation ambiguous. Task-quality estimates, instruction following, and preferences about continuation can produce similar choices. A successful avoidance attempt also requires the capability to carry it out. Consequently, counting such attempts cannot by itself distinguish stronger motivation from greater competence. Evaluation needs controlled behavioral comparisons that clarify what changes the model's choices, alongside a setting that reveals how those choices interact with resource constraints and other agents.

We examine an operational definition through three complementary experiments (Figure~\ref{fig:overview}). Refill decisions probe responses to resource depletion; matched decisions about another assistant provide a comparison; an arena enforces depletion while agents solve tasks and exchange resources. This progression connects controlled choices to survival performance without treating them as the same outcome. Our contributions are a framework for attributing self-preservation behavior, controlled comparisons for testing it, and an arena protocol with a proposed survival-based league leaderboard.
